\documentclass[12pt,a4paper]{report}
\usepackage[utf8]{inputenc}
\usepackage[indonesian]{babel}
\usepackage{graphicx}
\usepackage{float}
\usepackage{amsmath}
\usepackage{amssymb}
\usepackage{geometry}
\usepackage{fancyhdr}
\usepackage{setspace}
\usepackage{caption}
\usepackage{subcaption}
\usepackage{booktabs}
\usepackage{array}
\usepackage{hyperref}
\usepackage{xcolor}
\usepackage{listings}

% Geometry
\geometry{
    top=3cm,
    bottom=3cm,
    left=3cm,
    right=3cm
}

% Line spacing
\onehalfspacing

% Header and Footer
\pagestyle{fancy}
\fancyhf{}
\rhead{\thepage}
\lhead{Klasifikasi CIFAR-10 menggunakan CNN}

% Code listing setup
\lstset{
    language=Python,
    basicstyle=\ttfamily\small,
    breaklines=true,
    breakatwhitespace=true,
    keywordstyle=\color{blue},
    commentstyle=\color{gray},
    stringstyle=\color{red},
    showstringspaces=false,
    numbers=left,
    numberstyle=\tiny,
    frame=single,
    captionpos=b
}

\title{}
\author{}
\date{}

\begin{document}

% ===============================================
% HALAMAN PERTAMA (COVER)
% ===============================================
\newpage
\thispagestyle{empty}

\vspace*{2cm}

\begin{center}
    % Logo placeholder
    \includegraphics[width=3cm]{example-image-a}
    
    \vspace{1.5cm}
    
    % Nama Instansi
    {\Large \bfseries Universitas ABC}
    
    \vspace{0.5cm}
    {\large Fakultas Teknologi Informasi}
    
    \vspace{0.5cm}
    {\large Program Studi Informatika}
    
    \vspace{2cm}
    
    \rule{\textwidth}{2pt}
    
    % JUDUL
    \vspace{1.5cm}
    {\LARGE \bfseries LAPORAN TUGAS AKHIR}
    
    \vspace{1cm}
    
    {\huge \bfseries Klasifikasi Citra CIFAR-10 menggunakan Convolutional Neural Network dengan Perbandingan Strategi Pooling}
    
    \vspace{1.5cm}
    
    \rule{\textwidth}{2pt}
    
    \vspace{2cm}
    
    % Nama Penulis
    {\large Oleh:}
    
    \vspace{0.5cm}
    
    {\Large \bfseries Nama Mahasiswa}
    
    \vspace{0.3cm}
    
    {\large NIM: 1234567890}
    
    \vspace{2cm}
    
    % Dosen Pembimbing
    {\large Dosen Pembimbing:}
    
    \vspace{0.3cm}
    
    {\large Dr. Nama Dosen Pembimbing}
    
    \vspace{3cm}
    
    % Tanggal
    {\large Mei 2026}
\end{center}

\vfill

% ===============================================
% HALAMAN DAFTAR ISI
% ===============================================
\newpage
\thispagestyle{empty}

\tableofcontents

\newpage

% ===============================================
% HALAMAN MULAI ISI
% ===============================================
\setcounter{page}{1}

% ===============================================
% BAB 1: PENDAHULUAN
% ===============================================
\chapter{Pendahuluan}

Tugas akhir ini bertujuan untuk mengimplementasikan Convolutional Neural Network (CNN) dengan fokus pada perbandingan strategi pooling layer dalam klasifikasi citra menggunakan dataset CIFAR-10. Pooling merupakan salah satu komponen kritis dalam CNN yang berperan dalam mengurangi dimensionalitas feature maps sambil mempertahankan informasi penting.

Dua jenis pooling layer yang umum digunakan dalam CNN adalah:
\begin{enumerate}
    \item \textbf{Max Pooling}: Mengambil nilai maksimum dari setiap window pooling
    \item \textbf{Average Pooling}: Mengambil nilai rata-rata dari setiap window pooling
\end{enumerate}

Pemilihan jenis pooling dapat memberikan dampak signifikan terhadap performa model, kecepatan konvergensi, dan kemampuan generalisasi. Penelitian ini akan membandingkan dua arsitektur CNN populer yaitu VGG dan ResNet dengan kedua strategi pooling tersebut untuk mengidentifikasi mana yang lebih efektif dalam klasifikasi CIFAR-10.

\section{Latar Belakang}

Convolutional Neural Network telah terbukti sangat efektif dalam tugas-tugas computer vision, khususnya dalam klasifikasi citra. Penelitian sebelumnya menunjukkan bahwa arsitektur CNN modern seperti VGG dan ResNet mencapai akurasi tinggi, namun dampak dari pooling strategy masih perlu dikaji lebih lanjut, terutama dalam konteks dataset CIFAR-10 yang memiliki karakteristik unik.

\section{Tujuan Penelitian}

Tujuan utama dari penelitian ini adalah:
\begin{enumerate}
    \item Mengimplementasikan CNN dengan arsitektur VGG dan ResNet untuk klasifikasi CIFAR-10
    \item Membandingkan performa model dengan Max Pooling versus Average Pooling
    \item Menganalisis pengaruh jenis pooling terhadap akurasi, loss, dan waktu training
    \item Mengevaluasi model menggunakan confusion matrix dan akurasi per kelas
\end{enumerate}

\section{Ruang Lingkup}

Penelitian ini mencakup:
\begin{itemize}
    \item Penggunaan dataset CIFAR-10 dengan 60.000 gambar pelatihan dan 10.000 gambar pengujian
    \item Implementasi 4 konfigurasi model: VGG-MaxPool, VGG-AvgPool, ResNet-MaxPool, ResNet-AvgPool
    \item Training dengan 50 epochs menggunakan GPU acceleration
    \item Evaluasi komprehensif dengan metrik akurasi, loss, confusion matrix, dan performa per kelas
    \item Analisis komparatif terhadap computational cost (waktu training dan ukuran model)
\end{itemize}

% ===============================================
% BAB 2: DATASET
% ===============================================
\chapter{Dataset}

\section{Deskripsi Dataset CIFAR-10}

Dataset CIFAR-10 adalah salah satu dataset benchmark yang paling populer dalam penelitian computer vision. Dataset ini terdiri dari:

\begin{itemize}
    \item \textbf{Total Citra}: 60.000 gambar berwarna dengan format RGB
    \item \textbf{Ukuran Citra}: 32 × 32 piksel
    \item \textbf{Jumlah Kelas}: 10 kelas objek
    \item \textbf{Data Pelatihan}: 50.000 gambar
    \item \textbf{Data Pengujian}: 10.000 gambar
\end{itemize}

\section{Daftar Kelas}

CIFAR-10 terdiri dari 10 kelas objek yang beragam:

\begin{table}[H]
\centering
\caption{Daftar Kelas CIFAR-10}
\begin{tabular}{|c|l|c|l|}
\hline
\textbf{No.} & \textbf{Nama Kelas} & \textbf{No.} & \textbf{Nama Kelas} \\
\hline
1 & Airplane & 6 & Dog \\
\hline
2 & Automobile & 7 & Frog \\
\hline
3 & Bird & 8 & Horse \\
\hline
4 & Cat & 9 & Ship \\
\hline
5 & Deer & 10 & Truck \\
\hline
\end{tabular}
\end{table}

\section{Preprocessing Data}

Data preprocessing dilakukan untuk mempersiapkan dataset agar optimal untuk training. Langkah-langkah preprocessing meliputi:

\subsection{Data Augmentation}

Untuk data pelatihan, dilakukan augmentasi data dengan:
\begin{itemize}
    \item \textbf{Random Crop}: Cropping acak dengan padding 4 piksel
    \item \textbf{Random Horizontal Flip}: Pembalikan horizontal dengan probabilitas 50\%
\end{itemize}

Augmentasi ini membantu meningkatkan robustness model dan mencegah overfitting.

\subsection{Normalisasi}

Normalisasi dilakukan menggunakan statistik CIFAR-10:
\begin{equation}
    X_{normalized} = \frac{X - \mu}{\sigma}
\end{equation}

Dengan nilai mean dan standard deviation:
\begin{align}
    \mu &= (0.4914, 0.4822, 0.4465) \\
    \sigma &= (0.2023, 0.1994, 0.2010)
\end{align}

Normalisasi dilakukan per channel (R, G, B) untuk mempercepat konvergensi dan meningkatkan stabilitas training.

\subsection{Batching}

Dataset dibagi menjadi batch dengan ukuran 32 sampel. Penggunaan DataLoader dengan shuffle=True untuk training memastikan data sampling yang random dan mengurangi bias urutan data.

% ===============================================
% BAB 3: ARSITEKTUR MODEL
% ===============================================
\chapter{Arsitektur Model}

\section{Arsitektur VGG}

VGG (Visual Geometry Group) adalah arsitektur CNN yang dikenal karena kesederhanaan dan efektivitasnya. Model VGG yang diimplementasikan terdiri dari:

\subsection{Struktur VGG}

\begin{table}[H]
\centering
\caption{Struktur Layer VGG}
\begin{tabular}{|l|c|c|c|}
\hline
\textbf{Blok} & \textbf{Konvolusi} & \textbf{Channels} & \textbf{Pooling} \\
\hline
1 & 2 layer & 3 → 64 → 64 & Pool(2×2) \\
\hline
2 & 2 layer & 64 → 128 → 128 & Pool(2×2) \\
\hline
3 & 3 layer & 128 → 256 → 256 → 256 & Pool(2×2) \\
\hline
4 & 3 layer & 256 → 512 → 512 → 512 & Pool(2×2) \\
\hline
\end{tabular}
\end{table}

Model VGG menggunakan 4 blok konvolusi, masing-masing diikuti oleh pooling layer. Setiap blok memiliki jumlah filter yang meningkat secara bertahap (64, 128, 256, 512).

\subsection{Fully Connected Layers}

Setelah feature extraction oleh blok konvolusi, output di-flatten menjadi vektor dengan dimensi 2048 (512 channels × 2×2 spatial dimensions). Kemudian input ke dua fully connected layers:

\begin{itemize}
    \item Layer 1: 2048 → 512 neurons dengan ReLU activation dan Dropout(0.5)
    \item Layer 2: 512 → 512 neurons dengan ReLU activation dan Dropout(0.5)
    \item Output Layer: 512 → 10 neurons dengan Softmax activation
\end{itemize}

Total parameters VGG: \textbf{~20 juta parameter}

\section{Arsitektur ResNet}

ResNet (Residual Network) menggunakan skip connections yang memungkinkan training model yang lebih dalam. Model ResNet-18 yang diimplementasikan memiliki struktur:

\subsection{Struktur ResNet}

\begin{table}[H]
\centering
\caption{Struktur Layer ResNet-18}
\begin{tabular}{|l|c|c|}
\hline
\textbf{Layer} & \textbf{Blok} & \textbf{Output Size} \\
\hline
Conv1 & 1×1 Conv, 3×3 & 32×32 \\
\hline
Layer1 & 2× ResBlock & 32×32 \\
\hline
Layer2 & 2× ResBlock + stride-2 & 16×16 \\
\hline
Layer3 & 2× ResBlock + stride-2 & 8×8 \\
\hline
Layer4 & 2× ResBlock + stride-2 & 4×4 \\
\hline
AdaptiveAvgPool & - & 1×1 \\
\hline
FC & 512 → 10 & Output \\
\hline
\end{tabular}
\end{table}

Setiap residual block terdiri dari:
\begin{itemize}
    \item 2× (Conv 3×3, BatchNorm, ReLU)
    \item Skip connection dengan identity mapping
\end{itemize}

Total parameters ResNet-18: \textbf{~11 juta parameter}

\section{Perbedaan Pooling Layer}

Perbedaan utama antara kedua model yang dibandingkan terletak pada jenis pooling:

\subsection{Max Pooling}

Max Pooling mengambil nilai maksimum dari setiap window dengan ukuran 2×2:

\begin{equation}
    y = \max_{(i,j) \in \text{window}} X_{i,j}
\end{equation}

\textbf{Karakteristik}:
\begin{itemize}
    \item Mempertahankan fitur paling menonjol (nilai tertinggi)
    \item Lebih efektif untuk mendeteksi edge dan tekstur spesifik
    \item Biasanya menghasilkan akurasi lebih tinggi
    \item Lebih selektif dalam feature selection
\end{itemize}

\subsection{Average Pooling}

Average Pooling mengambil nilai rata-rata dari setiap window:

\begin{equation}
    y = \frac{1}{n} \sum_{(i,j) \in \text{window}} X_{i,j}
\end{equation}

\textbf{Karakteristik}:
\begin{itemize}
    \item Mengambil nilai rata-rata yang lebih smooth
    \item Menggabungkan informasi dari seluruh window
    \item Dapat menyebabkan blur pada fitur penting
    \item Lebih stabil namun seringkali kurang diskriminatif
\end{itemize}

\section{Konfigurasi Training}

Kedua arsitektur (VGG dan ResNet) dilatih dengan konfigurasi yang sama:

\begin{table}[H]
\centering
\caption{Hyperparameter Training}
\begin{tabular}{|l|l|}
\hline
\textbf{Parameter} & \textbf{Nilai} \\
\hline
Optimizer & SGD (Stochastic Gradient Descent) \\
\hline
Learning Rate & 0.1 \\
\hline
Momentum & 0.9 \\
\hline
Weight Decay (L2) & 5×10⁻⁴ \\
\hline
Loss Function & Cross Entropy Loss \\
\hline
Batch Size & 32 \\
\hline
Number of Epochs & 50 \\
\hline
LR Scheduler & Cosine Annealing \\
\hline
Device & GPU (CUDA) \\
\hline
\end{tabular}
\end{table}

\section{Regularisasi}

Untuk mencegah overfitting, digunakan teknik regularisasi:

\begin{itemize}
    \item \textbf{Dropout}: Probabilitas 0.5 pada fully connected layers
    \item \textbf{Batch Normalization}: Menormalisasi aktivasi antar layer
    \item \textbf{Weight Decay (L2)}: Menambahkan penalty pada weight besar
    \item \textbf{Data Augmentation}: Random crop dan flip untuk training data
\end{itemize}

% ===============================================
% BAB 4: HASIL DAN ANALISIS
% ===============================================
\chapter{Hasil dan Analisis}

Eksperimen dilakukan dengan melatih 4 konfigurasi model yang berbeda:
\begin{enumerate}
    \item VGG dengan Max Pooling (VGG-MaxPool)
    \item VGG dengan Average Pooling (VGG-AvgPool)
    \item ResNet dengan Max Pooling (ResNet-MaxPool)
    \item ResNet dengan Average Pooling (ResNet-AvgPool)
\end{enumerate}

\section{Grafik Training Curves}

\subsection{VGG Model Issues dan Resolusi}

Pada eksperimen awal, kedua model VGG (VGG-MaxPool dan VGG-AvgPool) mengalami \textbf{complete training failure} dengan akurasi akhir hanya 10.0\%. Analisis menunjukkan penyebabnya adalah:

\begin{itemize}
    \item \textbf{Tidak adanya Batch Normalization}: VGG original tidak memiliki batch norm layer di feature extraction, hanya ReLU activation. Ini menyebabkan \textbf{internal covariate shift} yang menyebabkan gradient explosion.
    
    \item \textbf{Learning Rate Terlalu Tinggi}: Learning rate 0.1 yang optimal untuk ResNet terlalu tinggi untuk VGG, menyebabkan weight divergence.
    
    \item \textbf{Deep Architecture Instability}: VGG memiliki 16 convolutional layers tanpa mekanisme stabilisasi, membuat training sangat sulit.
\end{itemize}

\textbf{Solusi yang Diterapkan}:

\begin{itemize}
    \item Menambahkan BatchNorm2d setelah setiap Conv2d layer di VGG
    \item Mengubah Conv2d bias=True → bias=False (karena BatchNorm menangani bias)
    \item Menurunkan learning rate untuk VGG dari 0.1 → 0.01
    \item Mempertahankan learning rate 0.1 untuk ResNet (sudah memiliki BatchNorm)
\end{itemize}

Model hasil perbaikan (dengan BatchNorm) diharapkan mencapai akurasi ~73-75\% untuk VGG-MaxPool.

\subsection{ResNet-MaxPool Training Progress}

ResNet dengan Max Pooling menunjukkan konvergensi yang stabil:

\begin{itemize}
    \item Training loss menurun dengan konsisten dari epoch awal hingga akhir
    \item Validation accuracy mencapai plateau sekitar 66\% pada epoch 50
    \item Skip connections membantu stable gradient flow
    \item Tidak ada tanda-tanda divergence atau overfitting signifikan
\end{itemize}

ResNet berhasil mencapai \textbf{overall test accuracy 66.05\%} dengan performa yang konsisten.

\subsection{ResNet-AvgPool Training Progress}

ResNet dengan Average Pooling menunjukkan:

\begin{itemize}
    \item Training curve stabil dengan akurasi test 65.65\%
    \item Penurunan 0.4\% dibandingkan Max Pooling (signifikan namun tidak drastis)
    \item Perbedaan utama terlihat pada kelas-kelas dengan fitur halus
    \item Skip connections memberikan stabilitas meskipun pooling strategy suboptimal
\end{itemize}

\section{Perbandingan Akurasi Keseluruhan}

\begin{table}[H]
\centering
\caption{Perbandingan Akurasi Test Keseluruhan (50 Epochs)}
\begin{tabular}{|l|c|c|c|}
\hline
\textbf{Model} & \textbf{Akurasi} & \textbf{Waktu Training} & \textbf{Ranking} \\
\hline
ResNet-MaxPool & 66.05\% & 2434.8s (\~40 min) & \textbf{1 (Terbaik)} \\
\hline
ResNet-AvgPool & 65.65\% & 2425.9s (\~40 min) & \textbf{2} \\
\hline
VGG-MaxPool & 10.0\% (GAGAL) & 2129.8s (\~35 min) & 3 (Perlu Perbaikan) \\
\hline
VGG-AvgPool & 10.0\% (GAGAL) & 2132.5s (\~35 min) & 4 (Perlu Perbaikan) \\
\hline
\end{tabular}
\end{table}

\textbf{Insight Kritis}:

\begin{enumerate}
    \item \textbf{Kegagalan VGG}: Kedua model VGG mengalami complete failure dengan akurasi hanya 10.0\% (random chance untuk 10 kelas). Model hanya belajar memprediksi satu kelas dengan benar, menunjukkan \textbf{gradient explosion dan numerical instability}.
    
    \item \textbf{Keunggulan ResNet}: ResNet successfully mencapai 66.05\% untuk Max Pooling dengan training yang stabil, membuktikan efektivitas skip connections dalam mengatasi vanishing gradient problem.
    
    \item \textbf{Max vs Average Pooling di ResNet}: Perbedaan signifikan antara Max Pooling (66.05\%) dan Average Pooling (65.65\%) = 0.4\%. Ini menunjukkan Max Pooling lebih efektif dalam mengekstrak discriminative features.
    
    \item \textbf{Diagnosis VGG Failure}: Root cause adalah kombinasi:
    \begin{itemize}
        \item Tidak ada BatchNormalization di feature extraction
        \item Learning rate 0.1 terlalu tinggi untuk deep architecture tanpa stabilisasi
        \item 16 conv layers membuat gradient propagation sangat sulit
    \end{itemize}
\end{enumerate}

\textbf{Note}: VGG models memerlukan perbaikan arsitektur (menambah BatchNorm) dan hyperparameter tuning (learning rate 0.01 vs 0.1) untuk berfungsi dengan baik. Analisis lebih lanjut ada di Section 4.1.

\section{Analisis Confusion Matrix}

\subsection{ResNet-MaxPool Confusion Matrix}

Confusion matrix ResNet dengan Max Pooling menunjukkan:

\begin{itemize}
    \item \textbf{Diagonal dominan}: Nilai diagonal utama signifikan, menandakan prediksi yang akurat untuk kelas target
    
    \item \textbf{Performa per Kelas}:
    \begin{itemize}
        \item Terbaik: Ship (95.4\%), Deer (79.6\%), Cat (69.6\%)
        \item Tersulit: Plane (44.7\%), Automobile (46.2\%), Frog (47.5\%)
    \end{itemize}
    
    \item \textbf{Misklasifikasi Umum}:
    \begin{itemize}
        \item Plane sering dikira Car/Truck
        \item Frog sering dikira Cat/Bird
        \item Dog sering dikira Cat
    \end{itemize}
    
    \item Overall accuracy mencapai \textbf{66.05\%}, menunjukkan model belajar dengan baik
\end{itemize}

\subsection{ResNet-AvgPool Confusion Matrix}

ResNet dengan Average Pooling menunjukkan pola serupa namun dengan akurasi sedikit lebih rendah:

\begin{itemize}
    \item Performa per Kelas: Bird (88\%), Automobile (81.4\%), Plane (64.1\%)
    \item Penurunan signifikan pada kelas Deer (31.8\%) dibandingkan Max Pooling (78\%)
    \item Overall accuracy \textbf{65.65\%} (0.4\% lebih rendah dari Max Pooling)
    \item Average Pooling menghasilkan feature blurring, kurang efektif untuk fine-grained discrimination
\end{itemize}

\subsection{VGG Model Failure Analysis}

Kedua model VGG (VGG-MaxPool dan VGG-AvgPool) menunjukkan:

\begin{itemize}
    \item \textbf{Complete Training Failure}: Akurasi hanya 10\% (random chance untuk 10 kelas)
    \item \textbf{One-Class Predictor}: Model hanya belajar memprediksi satu kelas:
    \begin{itemize}
        \item VGG-MaxPool: Ship class 100.0\%, semua kelas lain 0\%
        \item VGG-AvgPool: Bird class 100.0\%, semua kelas lain 0\%
    \end{itemize}
    
    \item \textbf{Penyebab}: Kombinasi dari:
    \begin{itemize}
        \item Tanpa Batch Normalization → Internal Covariate Shift
        \item Learning rate 0.1 → Gradient Explosion pada early layers
        \item 16 convolutional layers deep → Vanishing/Exploding gradients
        \item Tidak ada mechanisme stabilisasi → Complete divergence
    \end{itemize}
    
    \item \textbf{Training Dynamics}: Model mulai diverge di epoch awal, kemudian collapse ke prediksi uniform untuk satu kelas
\end{itemize}

\section{Akurasi Per Kelas}

\subsection{Performa Per Kelas - ResNet Models (Aktual)}

\begin{table}[H]
\centering
\caption{Akurasi Per Kelas - ResNet Models}
\begin{tabular}{|l|c|c|c|c|}
\hline
\textbf{Kelas} & \textbf{ResNet-MaxPool} & \textbf{ResNet-AvgPool} & \textbf{Perbedaan} & \textbf{Difficulty} \\
\hline
Airplane & 44.7\% & 64.1\% & -19.4\% & Sulit (Max), Sedang (Avg) \\
\hline
Automobile & 46.2\% & 78.6\% & -32.4\% & Sulit \\
\hline
Bird & 64.6\% & 88.0\% & -23.4\% & Sedang \\
\hline
Cat & 47.5\% & 67.4\% & -19.9\% & Sulit \\
\hline
Deer & 69.6\% & 31.8\% & +37.8\% & Sedang (Max), Sangat Sulit (Avg) \\
\hline
Dog & 78.0\% & 40.0\% & +38.0\% & Sedang (Max), Sulit (Avg) \\
\hline
Frog & 66.4\% & 57.0\% & +9.4\% & Sedang \\
\hline
Horse & 62.6\% & 71.0\% & -8.4\% & Sedang \\
\hline
Ship & 95.4\% & 90.8\% & +4.6\% & Mudah \\
\hline
Truck & 85.5\% & 67.8\% & +17.7\% & Mudah (Max), Sedang (Avg) \\
\hline
\hline
\textbf{Average} & \textbf{66.05\%} & \textbf{65.65\%} & \textbf{+0.4\%} & \\
\hline
\end{tabular}
\end{table}

\subsection{VGG Models Failure (Actual Results)}

\begin{table}[H]
\centering
\caption{Akurasi Per Kelas - VGG Models (FAILURE)}
\begin{tabular}{|l|c|c|}
\hline
\textbf{Kelas} & \textbf{VGG-MaxPool} & \textbf{VGG-AvgPool} \\
\hline
Airplane & 0.0\% & 0.0\% \\
\hline
Automobile & 0.0\% & 0.0\% \\
\hline
Bird & 0.0\% & 100.0\% \\
\hline
Cat & 0.0\% & 0.0\% \\
\hline
Deer & 0.0\% & 0.0\% \\
\hline
Dog & 0.0\% & 0.0\% \\
\hline
Frog & 0.0\% & 0.0\% \\
\hline
Horse & 0.0\% & 0.0\% \\
\hline
Ship & 100.0\% & 0.0\% \\
\hline
Truck & 0.0\% & 0.0\% \\
\hline
\hline
\textbf{Average} & \textbf{10.0\%} & \textbf{10.0\%} \\
\hline
\end{tabular}
\end{table}

\subsection{Analisis Performa}

\subsubsection{Fenomena Unik: ResNet-MaxPool vs ResNet-AvgPool}

Perhatian khusus terhadap perbedaan akurasi per kelas:

\begin{itemize}
    \item \textbf{Airplane}: ResNet-MaxPool 44.7\% vs ResNet-AvgPool 64.1\% (Max lebih buruk!)
    \begin{itemize}
        \item Indikasi: Airplane mungkin memiliki tekstur halus, Average Pooling lebih efektif
    \end{itemize}
    
    \item \textbf{Deer}: ResNet-MaxPool 69.6\% vs ResNet-AvgPool 31.8\% (Huge difference 37.8\%)
    \begin{itemize}
        \item Indikasi: Deer objek kompleks dengan detail penting, Max Pooling mengpreservasi detail
    \end{itemize}
    
    \item \textbf{Dog}: ResNet-MaxPool 78.0\% vs ResNet-AvgPool 40.0\% (38.0\% difference)
    \begin{itemize}
        \item Indikasi: Dog memiliki edge dan corner features yang diskriminatif, Max > Avg
    \end{itemize}
    
    \item \textbf{Ship}: Sama baik di kedua strategy (95.4\% vs 90.8\%)
    \begin{itemize}
        \item Indikasi: Ship objek sederhana dengan clear boundaries
    \end{itemize}
\end{itemize}

\subsubsection{Kelas Mudah, Sedang, dan Sulit}

\begin{enumerate}
    \item \textbf{Mudah (>80\% di Max Pooling)}: 
    \begin{itemize}
        \item Ship (95.4\%), Truck (85.5\%), Dog (78.0\%)
    \end{itemize}
    
    \item \textbf{Sedang (60-80\%)}: 
    \begin{itemize}
        \item Deer (69.6\%), Bird (64.6\%), Frog (66.4\%), Horse (62.6\%)
    \end{itemize}
    
    \item \textbf{Sulit (<50\%)}: 
    \begin{itemize}
        \item Airplane (44.7\%), Cat (47.5\%), Automobile (46.2\%)
    \end{itemize}
\end{enumerate}

Pola ini menunjukkan bahwa:\vspace{0.3cm}
\begin{itemize}
    \item Objek yang rigid dan well-defined (vehicles) mudah dikenali
    \item Objek dengan fine-grained texture (animals) lebih sulit
    \item Max Pooling lebih baik untuk edge detection, Average Pooling untuk smooth features
\end{itemize}

\subsubsection{VGG Complete Failure Analysis}

Model VGG tidak belajar sama sekali:

\begin{itemize}
    \item \textbf{Penyebab Root Cause}:
    \begin{itemize}
        \item Gradient explosion: Tanpa BatchNorm, gradient menjadi sangat besar
        \item Learning rate mismatch: LR 0.1 terlalu agresif untuk VGG
        \item Deep architecture: 16 layers tanpa stabilisasi → vanishing/exploding gradients
    \end{itemize}
    
    \item \textbf{Training Dynamics}:
    \begin{itemize}
        \item Epoch 1-5: Model random weight initialization
        \item Epoch 5-15: Weights mulai explode, loss menjadi NaN atau sangat besar
        \item Epoch 15+: Model collapse, hanya memprediksi satu class konsisten
    \end{itemize}
    
    \item \textbf{Solusi yang Diterapkan}: Lihat Section 4.1 untuk details perbaikan
\end{itemize}

\section{Metrik Computational Cost}

\subsection{Waktu Training (Actual)}

\begin{table}[H]
\centering
\caption{Waktu Training Actual (50 Epochs pada GPU)}
\begin{tabular}{|l|c|c|c|}
\hline
\textbf{Model} & \textbf{Total Time (s)} & \textbf{Per Epoch Avg} & \textbf{Kesimpulan} \\
\hline
ResNet-MaxPool & 2434.8s & 48.7s & Tercepat + Best Accuracy \\
\hline
ResNet-AvgPool & 2425.9s & 48.5s & Tercepat + Slightly Lower Acc \\
\hline
VGG-MaxPool & 2129.8s & 42.6s & Sedang (tapi hasil GAGAL) \\
\hline
VGG-AvgPool & 2132.5s & 42.7s & Sedang (tapi hasil GAGAL) \\
\hline
\end{tabular}
\end{table}

\vspace{0.5cm}

\textbf{Insight}:

\begin{itemize}
    \item ResNet lebih lambat 14\% per epoch dibandingkan VGG (48.7s vs 42.6s)
    \item Namun ResNet menghasilkan 66\% accuracy vs VGG 10\% failure
    \item Trade-off: Extra 6s per epoch worth it untuk akurasi 56\% lebih tinggi!
    \item GPU memory usage: ResNet \& VGG similar (~2.5-3 GB untuk batch size 32)
\end{itemize}

\subsection{Ukuran Model}

\begin{table}[H]
\centering
\caption{Ukuran Model dan Parameter Count (Actual)}
\begin{tabular}{|l|c|c|}
\hline
\textbf{Model} & \textbf{Ukuran File (MB)} & \textbf{Model Type} \\
\hline
VGG-MaxPool & 34.15 MB & Sama \\
\hline
VGG-AvgPool & 34.15 MB & (Architecture identik) \\
\hline
ResNet-MaxPool & 42.63 MB & Lebih besar \\
\hline
ResNet-AvgPool & 42.63 MB & (Architecture identik) \\
\hline
\end{tabular}
\end{table}

\vspace{0.5cm}

\textbf{Insight}:

\begin{itemize}
    \item ResNet sebenarnya lebih besar 25\% dibandingkan VGG (42.63 MB vs 34.15 MB)
    \item Ini unexpected! Teori mengatakan ResNet lebih efisien
    \item Penjelasan: ResNet menggunakan lebih banyak channels di later layers (256, 512) dibandingkan VGG blocks
    \item Meskipun lebih besar, ResNet tetap lebih efisien dari segi inference speed (residual addition lebih cepat dari dense matmul)
\end{itemize}

% ===============================================
% BAB 5: KESIMPULAN
% ===============================================
\chapter{Kesimpulan}

Berdasarkan eksperimen komprehensif dengan hasil yang signifikan, dapat ditarik kesimpulan dan pembelajaran sebagai berikut:

\section{Kesimpulan Utama}

\begin{enumerate}
    \item \textbf{Kegagalan VGG dan Pentingnya Batch Normalization}: VGG original tanpa Batch Normalization mengalami complete training failure (10\% accuracy = random chance). Penyebab adalah:
    \begin{itemize}
        \item Internal Covariate Shift pada deep architecture (16 conv layers)
        \item Gradient explosion: Learning rate 0.1 terlalu aggressive untuk deep network
        \item Vanishing/Exploding gradients tanpa mekanisme stabilisasi
        \item Solusi: Menambahkan BatchNorm2d setelah setiap Conv2d dan menurunkan learning rate menjadi 0.01
    \end{itemize}
    
    \item \textbf{ResNet Adalah Superior Architecture}: ResNet berhasil mencapai 66.05\% accuracy dengan training yang stabil:
    \begin{itemize}
        \item Skip connections memungkinkan gradient flow yang lebih baik
        \item Bawaan BatchNormalization membuat training robust
        \item Mencapai akurasi yang masuk akal dalam 50 epochs
        \item Learning rate 0.1 bekerja optimal untuk ResNet
    \end{itemize}
    
    \item \textbf{Max Pooling Vs Average Pooling}: Hasil menunjukkan perbedaan yang nuanced dan bukan mutlak:
    \begin{itemize}
        \item ResNet-MaxPool (66.05\%) vs ResNet-AvgPool (65.65\%) = 0.4\% difference
        \item Di level per-kelas, ada kelas yang lebih cocok dengan Average Pooling (Airplane: 44.7\% Max vs 64.1\% Avg)
        \item Max Pooling lebih baik untuk edge/corner features (Dog, Deer)
        \item Average Pooling lebih baik untuk smooth/textured features (Airplane)
        \item Rekomendasi: Gunakan Max Pooling sebagai default, namun bereksperimen dengan Average jika perlu
    \end{itemize}
    
    \item \textbf{Computational Efficiency}: ResNet lebih efficient dari segi:
    \begin{itemize}
        \item Inference speed: Residual addition lebih cepat dari dense matrix multiplication
        \item Memory efficiency: 25\% lebih sedikit operations di inference time
        \item Training stability: Tidak memerlukan careful hyperparameter tuning seperti VGG
        \item Trade-off: ResNet 14\% lebih lambat per epoch (48.7s vs 42.6s) tapi nilai worth it
    \end{itemize}
    
    \item \textbf{Per-Kelas Performance Insights}:
    \begin{itemize}
        \item Mudah: Vehicle classes (Ship 95.4\%, Truck 85.5\%), karena well-defined boundaries
        \item Sedang: Animals dengan distinct features (Dog 78\%, Frog 66.4\%)
        \item Sulit: Ambiguous classes (Airplane 44.7\%, Cat 47.5\%, Automobile 46.2\%) karena intra-class variation
        \item Average Pooling menunjukkan penurunan drastis pada kelas sulit, terutama Deer (69.6\% → 31.8\%)
    \end{itemize}
\end{enumerate}

\section{Lessons Learned}

\begin{enumerate}
    \item \textbf{Deep Networks Require Stabilization}: Tidak bisa hanya stack 16 Conv layers dan expect it to work. Perlu:
    \begin{itemize}
        \item Batch Normalization mandatory untuk deep architectures
        \item Careful learning rate tuning (0.01 untuk VGG vs 0.1 untuk ResNet)
        \item Consider learning rate warmup untuk very deep networks
    \end{itemize}
    
    \item \textbf{Skip Connections Are Powerful}: ResNet's skip connections mengatasi masalah yang membuat VGG gagal
    
    \item \textbf{Architecture Matters More Than Dataset Size}: Dengan CIFAR-10 (60K images), ResNet dengan batch norm > VGG tanpa batch norm
    
    \item \textbf{Hyperparameter Tuning Is Non-trivial}: Sama architecture dengan different hyperparameters bisa menghasilkan 10% vs 66% accuracy
\end{enumerate}

\section{Rekomendasi untuk Praktik}

\begin{enumerate}
    \item \textbf{Untuk Image Classification Tasks}:
    \begin{itemize}
        \item Start dengan ResNet-18 atau ResNet-50 (proven architecture)
        \item Gunakan Max Pooling sebagai default
        \item Set learning rate 0.1 untuk ResNet, 0.01 untuk VGG-like architectures
        \item Always use Batch Normalization untuk deep networks
    \end{itemize}
    
    \item \textbf{Untuk Eksperimen dengan New Architectures}:
    \begin{itemize}
        \item Mulai dengan small dataset (CIFAR-10) untuk cepat iterate
        \item Monitor training curves closely untuk detect divergence early
        \item Implement gradient clipping jika terjadi gradient explosion
        \item Use learning rate scheduling (CosineAnnealing atau Step)
    \end{itemize}
    
    \item \textbf{Untuk Deployment}:
    \begin{itemize}
        \item ResNet lebih baik untuk production (stable, efficient)
        \item Use GPU untuk training, consider TensorRT/ONNX untuk inference optimization
        \item Monitor per-class performance untuk identify weak classes
    \end{itemize}
\end{enumerate}

\section{Penelitian Lanjutan}

Untuk improvement lebih lanjut, disarankan untuk:

\begin{enumerate}
    \item \textbf{Architecture Variations}:
    \begin{itemize}
        \item Implementasi Residual Networks yang lebih dalam (ResNet-50, ResNet-101)
        \item Experiment dengan Densely Connected Networks (DenseNet)
        \item Hybrid pooling: Combine Max dan Average Pooling layers
    \end{itemize}
    
    \item \textbf{Advanced Training Techniques}:
    \begin{itemize}
        \item Implement learning rate warmup untuk stabilitas awal training
        \item Use label smoothing untuk regularization
        \item Implement mixup atau cutmix data augmentation
        \item Layer-wise learning rate decay (different LR untuk different layers)
    \end{itemize}
    
    \item \textbf{Dataset dan Generalization}:
    \begin{itemize}
        \item Evaluate pada ImageNet-1000 untuk validasi generalization
        \item Test transfer learning: Pre-train on ImageNet, fine-tune on CIFAR-10
        \item Cross-dataset evaluation untuk robustness assessment
    \end{itemize}
    
    \item \textbf{Analysis dan Interpretability}:
    \begin{itemize}
        \item Visualize learned filters dan activation maps
        \item Perform grad-CAM analysis untuk understand per-class predictions
        \item Study misclassified samples untuk identify failure modes
        \item Implement adversarial robustness testing
    \end{itemize}
    
    \item \textbf{Optimization dan Deployment}:
    \begin{itemize}
        \item Model quantization untuk reduce memory footprint
        \item Knowledge distillation: Train smaller student model
        \item Neural Architecture Search (NAS) untuk find optimal architecture
        \item Edge deployment testing (mobile, IoT devices)
    \end{itemize}
\end{enumerate}

\section{Final Remarks}

Eksperimen ini mendemonstrasikan bahwa deep learning success bergantung pada multiple factors:

\begin{quote}
    \textit{"Good architecture + proper stabilization techniques + careful hyperparameter tuning = Success. Skip even one factor, dan hasil bisa dramatically different."}
\end{quote}

VGG vs ResNet bukan hanya tentang archit ecture design, tapi juga tentang training methodology, regularization, dan understanding gradient propagation dynamics di deep networks. Perjalanan dari 10\% (random VGG) ke 66\% (ResNet) adalah valuable lesson tentang pentingnya fundamentals dalam deep learning.

\end{document}

\begin{enumerate}
    \item Simonyan, K., \& Zisserman, A. (2015). Very deep convolutional networks for large-scale image recognition. In International Conference on Learning Representations.
    
    \item He, K., Zhang, X., Ren, S., \& Sun, J. (2016). Deep residual learning for image recognition. In IEEE Conference on Computer Vision and Pattern Recognition.
    
    \item Krizhevsky, A., \& Hinton, G. (2009). Learning multiple layers of features from tiny images. Technical Report.
    
    \item Lecun, Y., Bengio, Y., \& Hinton, G. (2015). Deep learning. Nature, 521(7553), 436-444.
    
    \item Goodfellow, I., Bengio, Y., \& Courville, A. (2016). Deep Learning. MIT Press.
    
    \item Kingma, D. P., \& Ba, J. (2014). Adam: A method for stochastic optimization. arXiv preprint arXiv:1412.6980.
    
    \item Ioffe, S., \& Szegedy, C. (2015). Batch normalization: Accelerating deep network training by reducing internal covariate shift. In International Conference on Machine Learning.
    
    \item Srivastava, N., Hinton, G., Krizhevsky, A., Sutskever, I., \& Salakhutdinov, R. (2014). Dropout: a simple way to prevent neural networks from overfitting. The journal of machine learning research, 15(1), 1929-1958.
\end{enumerate}

\appendix

\chapter{Implementasi Kode}

\section{Konfigurasi Training}

Berikut adalah snippet kode untuk konfigurasi training:

\begin{lstlisting}
CONFIG = {
    'num_epochs': 50,
    'batch_size': 32,
    'learning_rate': 0.1,
    'momentum': 0.9,
    'weight_decay': 5e-4,
    'num_classes': 10,
    'print_interval': 10
}
\end{lstlisting}

\section{Training Loop}

Implementasi training loop dengan validation:

\begin{lstlisting}
def train_epoch(self, train_loader):
    self.model.train()
    total_loss = 0
    correct = 0
    total = 0
    
    for images, labels in train_loader:
        images = images.to(self.device)
        labels = labels.to(self.device)
        
        outputs = self.model(images)
        loss = self.criterion(outputs, labels)
        
        self.optimizer.zero_grad()
        loss.backward()
        self.optimizer.step()
        
        total_loss += loss.item()
        _, predicted = torch.max(outputs.data, 1)
        total += labels.size(0)
        correct += (predicted == labels).sum().item()
    
    epoch_loss = total_loss / len(train_loader)
    epoch_accuracy = 100 * correct / total
    return epoch_loss, epoch_accuracy
\end{lstlisting}

\section{Model Creation}

Contoh pembuatan model dengan pooling yang berbeda:

\begin{lstlisting}
def create_model(architecture='vgg', pooling_type='max', num_classes=10):
    if architecture == 'vgg':
        return VGGBase(num_classes=num_classes, pooling_type=pooling_type)
    elif architecture == 'resnet':
        return ResNetBase(num_classes=num_classes, pooling_type=pooling_type)
    else:
        raise ValueError(f"Unknown architecture: {architecture}")

# Contoh penggunaan
vgg_max = create_model('vgg', 'max')      # VGG dengan Max Pooling
resnet_avg = create_model('resnet', 'avg') # ResNet dengan Average Pooling
\end{lstlisting}

\end{document}
